\section*{Exercises on the Appendix}
\phantomsection
\addcontentsline{toc}{section}{Exercises on the Appendix}

\begin{enumerate}
\def\labelenumi{\arabic{enumi})}
\tightlist
\item
  Let \(\mathfrak{K}\) be a Hardy field such that, for every function \(f\in \mathfrak{K}\) not identically zero on a neighbourhood of \(+\infty\) there exists a \(\lambda >0\) such that
\end{enumerate}

\[
\frac {1}{e _ {m} (x ^ {\lambda})} \ll f (x) \ll e _ {m} (x ^ {\lambda})
\]

(m an integer independent of f).

\begin{enumerate}
\def\labelenumi{\alph{enumi})}
\tightlist
\item
  Let \(u_{1}, u_{2}, \ldots, u_{p}\) be p functions of the form \(u_{k} = \log |z_{k}|\) , where \(z_{k} \in K\) is not zero on a neighbourhood of \(+\infty\) . Show that for every function g (not zero on a neighbourhood of \(+\infty\) ) of the Hardy field \(\mathfrak{K}(u_{1}, u_{2}, \ldots, u_{p})\) obtained by adjoining the functions \(u_{k}\) ( \(1 \leqslant k \leqslant p\) ) to K, there exists a number \(\mu > 0\) such that
\end{enumerate}

\[
\frac {1}{e _ {m} (x ^ {\mu})} \ll g (x) \ll e _ {m} (x ^ {\mu})
\]

(reduce to the case where g is a polynomial in the \(u_{k}\) with coefficients in K and argue by induction on p, then, for p = 1, argue by induction on the degree of the polynomial g, proceeding as in lemma 2 of V, p.~248).

\begin{enumerate}
\def\labelenumi{\alph{enumi})}
\setcounter{enumi}{1}
\tightlist
\item
  Let \(u_{k}\) ( \(1 \leqslant k \leqslant p\) ) be \(p\) functions of the form \(u_{k} = \exp(z_{k})\) where \(z_{k} \in \mathcal{K}\) . Show that for every function \(g\) of the Hardy field \(\mathcal{K}(u_{1}, u_{2}, \ldots, u_{p})\) not identically zero on a neighbourhood of \(+\infty\) there exists a number \(\mu > 0\) such that
\end{enumerate}

\[
\frac {1}{e _ {m + 1} (x ^ {\mu})} \prec g (x) \prec e _ {m + 1} (x ^ {\mu})
\]

(similar method).

\begin{enumerate}
\def\labelenumi{\alph{enumi})}
\setcounter{enumi}{2}
\tightlist
\item
  Deduce that if \(f\) is an (H) function having a definition sequence with \(n\) terms, and not identically zero on a neighbourhood of \(+\infty\) , there exists a number \(\lambda > 0\) such that
\end{enumerate}

\[
\frac {1}{e _ {n} (x ^ {\lambda})} \ll f (x) \ll e _ {n} (x ^ {\lambda}).
\]

\begin{enumerate}
\def\labelenumi{\arabic{enumi})}
\setcounter{enumi}{1}
\tightlist
\item
  \begin{enumerate}
  \def\labelenumii{\alph{enumii})}
  \tightlist
  \item
    Show that every (H) function having a definition sequence of just one term is equivalent to a function of one of the forms \(x^{p}(\log x)^{q}\) , or \(x^{p}e^{g(x)}\) , where p and q are rational integers and g is a polynomial in x (method of exerc. 1).
  \end{enumerate}
\end{enumerate}

\begin{enumerate}
\def\labelenumi{\alph{enumi})}
\setcounter{enumi}{1}
\tightlist
\item
  Deduce from a) that every function in the Hardy field \(\mathbf{R}(x,u_{1},\ldots,u_{p})\) , where \(u_{1},\ldots,u_{p}\) are (H) functions having a single-term definition sequence, is equivalent to a function of the form \(x^{p}(\log x)^{q}e^{g(x)}\) where p and q are rational integers and g is a polynomial in x.
\end{enumerate}

\begin{enumerate}
\def\labelenumi{\arabic{enumi})}
\setcounter{enumi}{2}
\item
  Let f and g be two (H) functions such that f/g is of order 0 relative to \(l_{m}(x)\) ; show that if g is not of order 0 relative to \(l_{m}(x)\) one has \(f'/g' \sim f/g\) (compare \(\log |f|\) and \(\log |g|\) ).
\item
  Let \(\mathfrak{K}\) be a Hardy field such that, for every function \(f\in \mathfrak{K}\) not equivalent to a constant, and of order 0 relative to \(l_{m - 1}(x)\) , there exist a constant \(k\) and a rational integer \(r\) such that \(f(x)\sim k(l_m(x))^r\) .
\end{enumerate}

\begin{enumerate}
\def\labelenumi{\alph{enumi})}
\item
  Let z be any function in K, not identically zero on a neighbourhood of \(+\infty\) . Show that every function g in the Hardy field \(\mathcal{K}(\log|z|)\) not equivalent to a constant, and of order 0 relative to \(l_{m}(x)\) , is equivalent to a function of the form \(k(l_{m+1}(x))^{r}\) (k constant, r a rational integer). (First consider the case where g is a polynomial of degree p in \(\log|z|\) , with coefficients in K, and argue by induction on p, using exerc. 3; if g is a rational function of \(\log|z|\) , with coefficients in K, argue by induction on the degree of the numerator, using exerc. 3.)
\item
  Show that every function in \(\mathfrak{K}(l_{m+1}(x))\) is equivalent to a function of the form \(f(x)(l_{m+1}(x))^r\) where \(f \in \mathfrak{K}\) and \(r\) is a rational integer.
\item
  Let z be any function in K. Show that every function g of the Hardy field \(\mathfrak{K}(e^{z})\) of order 0 relative to \(l_{m}(x)\) is equivalent to a constant. (First consider the case where \(g = u^{q} \mathrm{P}(u)\) , where \(u = e^{z}\) , q is a rational integer, and \(\mathrm{P}(u)\) is a polynomial in u, of degree p, with coefficients in K; then argue by induction on p, using a) and exerc. 3. Pass from there to the general case by using exerc. 3.)
\item
  Extend the result of a) to the field \(\mathfrak{K}(u_{1}, u_{2}, \ldots, u_{s})\) , where the \(u_{k}\) are of the form \(e^{z_{k}}\) or \(\log |z_{k}|\) , the \(z_{k}\) being functions in K not identically zero on a neighbourhood of \(+\infty\) (argue by induction on s).
\end{enumerate}

\begin{enumerate}
\def\labelenumi{\arabic{enumi})}
\setcounter{enumi}{4}
\tightlist
\item
  \begin{enumerate}
  \def\labelenumii{\alph{enumii})}
  \tightlist
  \item
    Deduce from exerc. 4 that if f is an (H) function having a definition sequence of n terms, not equivalent to a constant, and of order 0 relative to \(l_{n-1}(x)\) , there exist a constant k and a rational integer r such that \(f(x) \sim k(l_n(x))^r\) .
  \end{enumerate}
\end{enumerate}

\begin{enumerate}
\def\labelenumi{\alph{enumi})}
\setcounter{enumi}{1}
\tightlist
\item
  Deduce from exerc. 4 that if \(f\) is any (H) function there exists an integer \(n\) such that the logarithmic criterion of order \(n\) will determine whether the integral \(\int_{a}^{+\infty} f(t) dt\) converges or is infinite.
\end{enumerate}

\begin{enumerate}
\def\labelenumi{\arabic{enumi})}
\setcounter{enumi}{5}
\tightlist
\item
  Compare among themselves the functions \(e_p\left((l_q(x))^{\mu}\right)\) according to the values of the integers \(p\) and \(q\) and the real number \(\mu\) , assumed \(> 0\) and \(\neq 1\) .
\end{enumerate}

7) Let \(f\) be an (H) function having a definition sequence of \(n\) terms. Show that if one has \(e_p((l_{q+1}(x))^{\beta}) \ll f(x) \ll e_p((l_q(x))^{\alpha})\) (resp. \(e_p((l_q(x))^{\beta}) \ll f(x) \ll e_{p+1}((l_q(x))^{\alpha}))\) )

for every \(\beta > 0\) and every \(\alpha\) such that \(0 < \alpha < 1\) , then necessarily \(p + q + 1 < n\) (argue by induction on \(n\) , using methods similar to those of exerc. 1 (V, p.~263) and 4 (V, p.~264)).

\begin{enumerate}
\def\labelenumi{\arabic{enumi})}
\setcounter{enumi}{7}
\tightlist
\item
  \begin{enumerate}
  \def\labelenumii{\alph{enumii})}
  \tightlist
  \item
    Let \((f_{n})\) be a sequence of increasing continuous functions belonging to \(\mathcal{H}(\mathfrak{F},\mathbf{R})\) and such that \(f_{n} \ll f_{n+1}\) for every n.~Show that there exists a function f that is continuous, increasing, belongs to \(\mathcal{H}(\mathfrak{F},\mathbf{R})\) , and is such that \(f \gg f_{n}\) for every n (reduce to the case where \(f_{n} \leqslant f_{n+1}\) and define f so that \(f_{n}(x) \leqslant f(x) \leqslant f_{n+1}(x)\) for \(n \leqslant x \leqslant n+1\) ).
  \end{enumerate}
\end{enumerate}

\begin{enumerate}
\def\labelenumi{\alph{enumi})}
\setcounter{enumi}{1}
\item
  Let \((f_n)\) be a sequence of increasing continuous functions belonging to \(\mathcal{H}(\mathfrak{F},\mathbf{R})\) and such that \(1\ll f_{n + 1}\ll f_n\) for every \(n\) . Show that there exists a function \(f\) that is continuous, increasing, belongs to \(\mathcal{H}(\mathfrak{F},\mathbf{R})\) , and is such that \(1\ll f\ll f_n\) for every \(n\) (reduce to the case where \(f_{n + 1}\leqslant f_n\) and show that one can define an increasing sequence \((x_{n})\) of real numbers and a continuous increasing function \(f\) such that \(f_{n + 1}(x)\leqslant f(x)\leqslant f_n(x)\) for \(x_{n}\leqslant x\leqslant x_{n + 1}\) ).
\item
  Let \((f_n), (g_n)\) be two sequences of continuous increasing functions belonging to \(\mathcal{H}(\mathfrak{F}, \mathbf{R})\) such that \(f_n \ll f_{n+1}\) , \(g_m \gg g_{m+1}\) and \(f_n \ll g_m\) for all \(m\) and \(n\) ; show that there exists a function \(h\) , continuous, increasing, belonging to \(\mathcal{H}(\mathfrak{F}, \mathbf{R})\) , such that \(f_n \ll h \ll g_m\) for all \(m\) and \(n\) (similar method).
\end{enumerate}

In particular, show that there exists a continuous decreasing function \(f\) in \(\mathcal{H}(\mathfrak{F},\mathbf{R})\) such that no logarithmic criterion allows one to determine whether the integral \(\int_{a}^{+\infty}f(t)dt\) converges or is infinite (``theorems of Du Bois-Reymond'').

\begin{enumerate}
\def\labelenumi{\arabic{enumi})}
\setcounter{enumi}{8}
\item
  Show that the series \(\sum_{n=1}^{\infty} \frac{e_n(x)}{e_n(n)}\) converges uniformly on every compact subset of \(\mathbf{R}\) , and that its sum \(f(x)\) is such that \(f(x) \gg e_n(x)\) for every integer \(n\) .
\item
  Show that there exists an increasing function \(f\) , defined, continuous and \(> 0\) for \(x \geqslant 0\) , such that \(f(2x) = 2^{f(x)}\) for all \(x \geqslant 0\) ; deduce from this that \(f(x) \gg e_n(x)\) for every integer \(n\) .
\item
  Let \(f\) be an increasing function, continuous and \(\geqslant 0\) , defined for \(x \geqslant 0\) , and such that \(f(x) \gg e_n(x)\) for every integer \(n\) ; show that, if \(g\) is the inverse function of \(f\) , one has \(1 \ll g(x) \ll l_n(x)\) for every integer \(n\) .
\end{enumerate}

12) For every integer \(n > 0\) let \(n = \sum_{k=0}^{\infty} \varepsilon_k 2^k\) be the dyadic expansion of \(n\) ( \(\varepsilon_k\) is an integer, zero for all but a finite number of indices, \(0 \leqslant \varepsilon_k \leqslant 1\) ). Put

\[
\alpha (n) = \sum_ {k = 0} ^ {\infty} \varepsilon_ {k}, \quad \mathrm{A} _ {1} (n) = \sum_ {j = 1} ^ {n} \alpha (j), \quad \mathrm{A} _ {2} (n) = \sum_ {j = 1} ^ {n} \alpha (\alpha (j)).
\]

\begin{enumerate}
\def\labelenumi{\alph{enumi})}
\tightlist
\item
  Prove the formula \(\mathrm{A}_1(n) = \sum_{k=0}^{\infty}(k+2)\varepsilon_k 2^{k-1}\) . Deduce the formula
\end{enumerate}

\[
\mathrm{A} _ {1} (n) = \frac {n}{2} \frac {\log n}{\log 2} + o (n \log \log n)
\]

as \(n\) grows indefinitely (decompose the sum which expresses \(\mathrm{A}_1(n)\) into two parts, the index \(k\) varying from 0 to \(\mu(n)\) in the first sum, from \(\mu(n)\) to \(n_1\) in the second, where \(n_1\) is the largest of the numbers \(k\) such that \(\varepsilon_k \neq 0\) and \(\mu(n)\) is chosen suitably; bound the first sum using prop. 6 of V, p.~239, and then estimate the difference between the second sum and \(n_1n / 2\) .

\begin{enumerate}
\def\labelenumi{\alph{enumi})}
\setcounter{enumi}{1}
\tightlist
\item
  Establish the relation
\end{enumerate}

\[
\sum_ {j = 0} ^ {2 ^ {m} - 1} f (\alpha (j)) = \sum_ {k = 0} ^ {m} \binom {m} {k} f (k)
\]

\(f\) being an arbitrary function defined on \(\mathbf{N}\) (for a given \(k\) consider the number of \(j < 2^{m}\) such that \(\alpha(j) = k\) ).

\begin{enumerate}
\def\labelenumi{\alph{enumi})}
\setcounter{enumi}{2}
\tightlist
\item
  For \(m = 2^r - 1\) show that \(\alpha(m - j) + \alpha(j) = r\) . Deduce from this and from \(b\) ) that
\end{enumerate}

\[
\mathrm{A} _ {2} (2 ^ {m}) = r 2 ^ {m - 1} \sim \frac {2 ^ {m} \log \log 2 ^ {m}}{2 \log 2}.
\]

\begin{enumerate}
\def\labelenumi{\alph{enumi})}
\setcounter{enumi}{3}
\tightlist
\item
  For \(m = 2^{r}\) show that
\end{enumerate}

\[
\mathrm{A} _ {2} (2 ^ {m}) = 2 \mathrm{A} _ {2} (2 ^ {m - 1}) + 2 ^ {m - 1} - \sum_ {j = 0} ^ {m - 1} \binom {m - 1} {j} \Delta (j + 1)
\]

where we have put \(\alpha(n+1)-\alpha(n)=1-\Delta(n+1)\) for all \(n\) (use \(b\) ) and the relation \(\alpha(2^{k}+r)=\alpha(r)+1\) for \(r<2^{k}\) ). Show that \(\sum_{j=0}^{m-1}\binom{m-1}{j}\Delta(j+1)\leqslant r^{m-2}\) (remark that \(\Delta(j+1)=0\) if \(j\) is even, and \(\Delta(j+1)\leqslant\log j/\log 2\) for all \(j\) ). Deduce, with the help of \(c\) ), that

\[
\mathrm{A} _ {2} (2 ^ {m}) \geqslant \frac {3}{2} r 2 ^ {m - 1} \geqslant \frac {3}{2} \frac {2 ^ {m} \log \log 2 ^ {m}}{2 \log 2}.
\]

Conclude from \(c\) ) and this relation that \(\mathrm{A}_2(n)\) is not equivalent to any (H) function as \(n\) tends to \(+\infty\) .

\begin{enumerate}
\def\labelenumi{\arabic{enumi})}
\setcounter{enumi}{12}
\tightlist
\item
  Let \(f\) be an (H) function such that \(1 \prec f(x) \prec x\) . Show that in the Taylor formula
\end{enumerate}

\[
\begin{array}{l} f (x + f (x)) = f (x) + f (x) f ^ {\prime} (x) + \frac {1}{2 !} (f (x)) ^ {2} f ^ {\prime \prime} (x) + \dots + \\ \qquad + \frac {1}{n !} (f (x)) ^ {n} f ^ {(n)} (x) + \frac {1}{(n + 1) !} (f (x)) ^ {n + 1} f ^ {(n + 1)} (x + \theta f (x)) \end{array}
\]

with \(0 < \theta < 1\) (I, p.~42, exerc. 9) each term is negligible relative to the previous one. If f is of order \textless{} 1 relative to x, the last term in this sum tends to 0 with 1/x once n is sufficiently large.

\begin{enumerate}
\def\labelenumi{\arabic{enumi})}
\setcounter{enumi}{13}
\tightlist
\item
  Deduce from exerc. 13 that if \(f\) is an (H) function such that \(\log f(e^x)\) is of order \(< 1\) relative to \(x\) , the function \(f(xf(x))\) is equivalent to a function of the form \(e^{g(x)}\) , where \(g\) is a rational function of \(x\) , \(\log f(x)\) , and of a certain number of derivatives of this last function.
\end{enumerate}

15) Let \(f\) be an (H) function, convex on a neighbourhood of \(+\infty\) , such that \(f(x) \gg x\) . For every \(\alpha > 0\) let \(x_{\alpha}\) be the point such that \(f'(x_{\alpha}) = (1 + \alpha)f'(x)\) .

\begin{enumerate}
\def\labelenumi{\alph{enumi})}
\tightlist
\item
  If \(f\) is of order \(+\infty\) relative to \(x\) show that
\end{enumerate}

\[
x _ {\alpha} - x \sim \log (1 + \alpha) \frac {f (x)}{f ^ {\prime} (x)}
\]

(use prop. 2 (V, p.~251) and 4 (V, p.~255), applying the Taylor formula to \(\log f'(x)\) ). Show that \(\left(f(x_{\alpha}) - (x_{\alpha} - x)f'(x_{\alpha})\right)/f(x)\) tends to \((1 + \alpha)\left(1 - \log(1 + \alpha)\right)\) as \(x\) tends to \(+\infty\) .

\begin{enumerate}
\def\labelenumi{\alph{enumi})}
\setcounter{enumi}{1}
\tightlist
\item
  If \(f\) is of order \(r > 1\) relative to \(x\) show that
\end{enumerate}

\[
x _ {\alpha} \sim (1 + \alpha) ^ {1 / (r - 1)} x
\]

(remark that if \(f_{1}\) is of order 0 relative to \(x\) one has \(f_{1}(kx) \sim f_{1}(x)\) for every constant \(k\) , by prop. 4 of V, p.~255). Deduce that \(\left(f(x_{\alpha}) - (x_{\alpha} - x)f'(x_{\alpha})\right) / f(x)\) tends to

\[
r (1 + \alpha) - (r - 1) (1 + \alpha) ^ {r / (r - 1)}
\]

as \(x\) tends to \(+\infty\) .

\begin{enumerate}
\def\labelenumi{\alph{enumi})}
\setcounter{enumi}{2}
\tightlist
\item
  Suppose finally that \(f\) is of order 1 relative to \(x\) . Putting \(f(x) = xf_1(x)\) show that
\end{enumerate}

\[
x _ {\alpha} - x \sim \log (1 + \alpha) \frac {f _ {1} (x)}{f _ {1} ^ {\prime} (x)}
\]

(consider the inverse function of \(f_{1}\) , which is of order \(+\infty\) relative to \(x\) ). Deduce that \(\left(f(x_{\alpha}) - (x_{\alpha} - x)f'(x_{\alpha})\right) / f(x)\) tends to \(-\infty\) as \(x\) tends to \(+\infty\) .

Similarly let \(x_{\alpha}'\) be the point such that \(f'(x_{\alpha}') = (1 - \alpha)f'(x)\) (for \(0 < \alpha < 1\) ). State the formulae analogous to the previous ones that express the principal part of \(x - x_{\alpha}'\) and the limit of

\[
\left(f (x _ {\alpha} ^ {\prime}) + (x - x _ {\alpha} ^ {\prime}) f ^ {\prime} (x _ {\alpha} ^ {\prime})\right) / f (x)
\]

as \(x\) tends to \(+\infty\) .

Conclude from these results that \(f\) is a regularly convex function on a neighbourhood of \(+\infty\) (V, p.~260, exerc. 5).

\makeatletter
\gdef\@chapapp{\chaptername}
\makeatother
\renewcommand{\thechapter}{\Roman{chapter}}
\renewcommand{\thesection}{\S\arabic{section}}
\setcounter{chapter}{5}
